\section{The invariant measure obstruction}
\label{sec:obstruction}

We now exclude the quadratic function produced by
Proposition~\ref{prop:limit} for the graph-defined bundle $G$.
The two rulings and the three graph identities force a nonzero squared seminorm on $\OO_{\PP^1}(-b)$
that is invariant under the lifts of $A$ and $B$.
We then turn that seminorm into a positive measure of mass $2\pi b$
and show that the two transformations cannot preserve such a measure.

\subsection{From the two rulings to one invariant seminorm}

\begin{lemma}\label{lem:ruling-reduction}
Let $G$ be the bundle of Proposition~\ref{prop:construction}, and put
$F=G^\vee$. Suppose that $q:F\to[0,\infty)$ is finite,
plurisubharmonic, not identically zero, and a Hermitian quadratic form
on every fiber. Then there is a finite, nonnegative, nonzero
plurisubharmonic function $a_\bullet$ on $\OO_{\PP^1}(-b)$, a Hermitian
positive semidefinite quadratic on each fiber, such that
\[
 a_\bullet(Mx,M^{\otimes b}v)=a_\bullet(x,v)
 \qquad(M=A,B)
\]
for every $x\in\PP^1$ and $v\in\OO(-b)_x$.
\end{lemma}

\begin{proof}
Write
$q(x,v)=H_x(v,v)$; nonnegativity says that every $H_x$ is positive
semidefinite.  Pull $q$ back by the line inclusions
\[
 \iota_i:L_i^{-1}=p_i^*\OO_{\PP^1}(-b)\longrightarrow F,
 \qquad i=1,2.
\]
For fixed $x_i\in\PP^1$ and $v\in\OO(-b)_{x_i}$, the function obtained
by holding $v$ fixed while varying the other base coordinate is finite
and subharmonic on the complementary $\PP^1$.  It is constant by the
maximum principle.  Consequently $q\circ\iota_i$ descends to a finite
nonnegative squared seminorm $a_i$ on $\OO_{\PP^1}(-b)$ over the
$i$-th factor.  Fixing any point in the other factor realizes $a_i$ as a
holomorphic pullback of $q$, so $a_i$ is plurisubharmonic on the total
line bundle.

On $\Gamma_M$, the identity $\iota_2\alpha=\iota_1$ from
Proposition~\ref{prop:construction} and the definition of $\alpha$ give
\begin{equation}
 \label{eq:seminorm-graph-identities}
 a_1(x,v)=a_2(Mx,M^{\otimes b}v),
 \qquad M\in\{\Id,A,B\}.
\end{equation}
The case $M=\Id$ identifies $a_1$ and $a_2$ as a single function
$a_\bullet$ on $\OO(-b)$.  This function is not identically zero.
Indeed, if it were zero, then for every $x\notin D$, the form $H_x$
would vanish on each of two spanning lines.  The inequality
\[
 |H_x(v,w)|^2\leq H_x(v,v)H_x(w,w)
\]
would then force $H_x=0$.  Hence $q$ would vanish on the full-measure
complement of the inverse image of $D$ in $F$.  Nonnegativity and the
submean inequality on small total-space coordinate balls would give $q=0$
everywhere, a contradiction.

The cases $M=A,B$ of \eqref{eq:seminorm-graph-identities} give the
asserted invariances under the holomorphic lifts $A^{\otimes b}$ and
$B^{\otimes b}$.
\end{proof}

\subsection{Logarithmic weights and their measure}

The seminorm in Lemma~\ref{lem:ruling-reduction} may vanish in some
fibers. The next lemma shows that its logarithmic weights still define
a positive measure with the mass dictated by the negative line bundle.
This is the rank-one squared-norm/logarithmic-norm principle from
singular Hermitian metric theory
\cite[Section~3]{BerndtssonPaun2008} and
\cite[Definition~1.2 and Section~2]{Raufi2015}.
We prove local integrability at zeros before taking distributional
curvature, so all possible atoms remain in the mass calculation.

\begin{lemma}[Logarithmic curvature of a squared seminorm]
\label{lem:seminorm-curvature}
Let $L$ be a holomorphic line bundle on a connected Riemann surface $C$.
Suppose that $a_\bullet:L\to[0,\infty)$ is finite, plurisubharmonic, and
not identically zero, and that its restriction to every fiber is a
Hermitian positive semidefinite quadratic form.  In a nonvanishing
holomorphic frame $e$ on $U\subset C$, write
\[
 a_\bullet(z,w e(z))=u_e(z)|w|^2.
\]
Then $\log u_e$, with the value $-\infty$ allowed, is subharmonic and
locally integrable.  The local measures $\ddc\log u_e$ define a positive
Radon measure $\mu$ on $C$.

If a holomorphic bundle automorphism $\widetilde T:L\to L$ covers an
automorphism $T:C\to C$ and satisfies
$a_\bullet\circ\widetilde T=a_\bullet$, then $T^*\mu=\mu$.
If $C=\PP^1$ and $L=\OO_{\PP^1}(-b)$ with $b\geq1$, then
\begin{equation}
 \label{eq:seminorm-curvature-mass}
 \mu(\PP^1)=2\pi b.
\end{equation}
\end{lemma}

Where $u_e$ is smooth and strictly positive, $\ddc\log u_e$ is the
negative of the Chern curvature of the corresponding metric in our
convention.
This accounts for the positive mass when $L=\OO(-b)$.

\begin{proof}
Fix a coordinate disk $U$ and a frame $e$, and abbreviate $u_e$ to $u$.
The function $u$ is finite, nonnegative, subharmonic, and locally bounded,
by restriction to the holomorphic section $e$.  Choose a nonnegative
smooth radial function $\rho$ supported in the unit disk, with integral
one, and put
\[
 \rho_\delta(z)=\delta^{-2}\rho(z/\delta),
 \qquad u_\delta=u*\rho_\delta+\delta.
\]
These functions are defined on the interior of $U$ for sufficiently
small $\delta>0$.  Averaging the holomorphic base translations
$(z,w)\mapsto(z-\zeta,w)$ shows that
$(u*\rho_\delta)(z)|w|^2$ is plurisubharmonic.  Adding
$\delta|w|^2$ preserves this property, so $u_\delta>0$ and
$u_\delta(z)|w|^2$ is smooth and plurisubharmonic.  Its Levi matrix at
$w=1$ is
\[
 \begin{pmatrix}
  (u_\delta)_{z\bar z} & (u_\delta)_z\\
  (u_\delta)_{\bar z} & u_\delta
 \end{pmatrix}.
\]
Its nonnegative determinant gives
\[
 (\log u_\delta)_{z\bar z}
 =\frac{u_\delta(u_\delta)_{z\bar z}
             -|(u_\delta)_z|^2}{u_\delta^2}\geq0.
\]
Thus $f_\delta:=\log u_\delta$ is subharmonic.

The radial convolutions are nonnegative weighted averages of circular
means of $u$.  These means increase with their radii, and converge to
$u$ as the radii tend to zero.  Consequently
\[
 u_\delta\downarrow u,
 \qquad f_\delta\downarrow\log u
 \quad\text{as }\delta\downarrow0.
\]
We establish local integrability before taking distributional
derivatives of the limit.  Suppose that $u(z_0)>0$, and choose $R>0$
with $\overline{B(z_0,R)}\subset U$.  For all sufficiently small
$\delta$, local boundedness supplies a constant $C_0\geq0$ such that
$f_\delta\leq C_0$ on this disk.  Moreover
$f_\delta(z_0)\geq\log u(z_0)$.  Integrating the circular submean
inequality over the radius gives
\[
 \int_{B(z_0,R)} f_\delta\,dA
 \geq \pi R^2\log u(z_0).
\]
Writing $f^-:=\max\{-f,0\}$, we obtain the explicit bound
\begin{equation}
 \label{eq:log-negative-part}
 \int_{B(z_0,R)} f_\delta^-\,dA
 \leq \pi R^2\bigl(C_0-\log u(z_0)\bigr).
\end{equation}
The negative parts increase as $\delta$ decreases, whereas the positive
parts are bounded by $C_0$.  Monotone convergence therefore proves
$\log u\in L^1(B(z_0,R))$.  It also gives
$f_\delta\to\log u$ in $L^1$ on this disk.

This conclusion propagates throughout $C$.  Indeed, let $\Omega$
consist of the points near which a local logarithmic weight is
integrable.  A frame change $e'=g e$ gives
\begin{equation}
 \label{eq:log-frame-change}
 \log u_{e'}=\log u_e+\log|g|^2,
\end{equation}
so $\Omega$ is independent of the chosen frame.  It is open and
nonempty: nontriviality of $a_\bullet$ supplies a point with $u>0$,
to which \eqref{eq:log-negative-part} applies.  If $x$ is in the
boundary of $\Omega$, choose a coordinate disk containing
$\overline{B(x,4r)}$.  There is a small open disk contained in
$\Omega\cap B(x,r)$.  On that disk integrability of $\log u$ implies
$u>0$ almost everywhere, so it contains a point $z_0$ with $u(z_0)>0$.
The disk $B(z_0,2r)$ has closure in the original chart and contains a
neighborhood of $x$.  Applying \eqref{eq:log-negative-part} on this
larger disk proves $x\in\Omega$.  Thus $\Omega$ is also closed, and
connectedness gives $\Omega=C$.

On every relatively compact disk, $C_0-f_\delta$ increases to the
integrable function $C_0-\log u$.  Hence the local convergence is in
$L^1$, and the decreasing limit $\log u$ is subharmonic.  In particular,
$\ddc\log u$ is a positive distribution and thus a positive Radon
measure.  Since $g$ in \eqref{eq:log-frame-change} is holomorphic and
nowhere zero, $\log|g|^2$ is harmonic.  The local measures therefore
glue to $\mu$.

For the invariance assertion, choose frames $e$ on $U$ and $e'$ on $V$
with $T(U)\subset V$, and write
\[
 \widetilde T(e(z))=c(z)e'(Tz),
\]
where $c$ is holomorphic and nowhere zero.  The isometry condition gives
\[
 u_e(z)=|c(z)|^2u_{e'}(Tz),
 \qquad
 \log u_e=\log|c|^2+(\log u_{e'})\circ T.
\]
Taking distributional $\ddc$, which commutes with pullback by the
holomorphic diffeomorphism $T$, gives $\mu=T^*\mu$.  Here pullback is
that of a top-degree current: for a test function $\varphi$,
\begin{equation}
 \label{eq:measure-pullback}
 \int_C\varphi\,d(T^*\mu)
   =\int_C(\varphi\circ T^{-1})\,d\mu.
\end{equation}
Equivalently, $(T^*\mu)(E)=\mu(T(E))$ for Borel sets $E$.  This is
ordinary measure invariance under the bijection $T$; in particular,
$T^*\delta_p=\delta_{T^{-1}p}$ with unchanged atom mass.

Finally suppose $C=\PP^1$ and $L=\OO(-b)$.  Let $h_0$ be the smooth
Hermitian metric on $\OO(-b)$ induced by the standard norm on $\CC^2$.
In the tautological affine frame $e(z)=(1,z)^{\otimes b}$ its weight is
\[
 u_0(z)=h_0(e(z),e(z))=(1+|z|^2)^b.
\]
The local functions $\log u_e-\log h_0(e,e)$ define a single globally
integrable function $\Phi$ on $\PP^1$, by the common frame law.
If $\mu_0=\ddc\log h_0(e,e)$, then
\[
 \mu-\mu_0=\ddc\Phi,
 \qquad
 \int_{\PP^1}(\mu-\mu_0)
   =\langle\ddc\Phi,1\rangle=0.
\]
The measure $\mu$ is finite because it is locally finite on the compact
sphere.  With the convention
$\ddc f=i f_{z\bar z}\,dz\wedge d\bar z$, the smooth reference measure
has mass
\[
 \int_{\PP^1}\mu_0
 =\int_{\CC}\frac{b\,i\,dz\wedge d\bar z}{(1+|z|^2)^2}
 =4\pi b\int_0^\infty\frac{r\,dr}{(1+r^2)^2}
 =2\pi b.
\]
This proves \eqref{eq:seminorm-curvature-mass}.  Zeros of $u_e$ and any
atoms of $\mu$ are included in the distributional calculation; no
points have been excised.
\end{proof}

\subsection{The incompatible invariances}

We now use the special dynamics of the two matrices. Concentration
of invariant measures on proper projective subsets is a classical
mechanism; see \cite[Lemma~2 and the proof of Lemma~3]{Furstenberg1976}.
The following elementary argument gives precisely the rank-two
statement needed here.

\begin{lemma}
\label{lem:no-invariant-measure}
The projective transformations induced by the matrices $A$ and $B$ in
\eqref{eq:matrices} preserve no nonzero finite positive
Borel measure on $\PP^1$.
\end{lemma}

\begin{proof}
Suppose that $\mu$ is such a measure.  It is invariant under $T=AB$.
The two eigenvalues of $AB$ have distinct moduli, so after a projective
change of coordinate $T$ acts as $z\mapsto\lambda z$ with
$R:=|\lambda|>1$.  The Borel sets
\[
 C_n=\{z\in\CC^*:R^n\leq|z|<R^{n+1}\},\qquad n\in\ZZ,
\]
are pairwise disjoint, cover $\CC^*$, and satisfy $T(C_n)=C_{n+1}$.
Each has the same measure.  Finiteness and positivity imply that this
common measure is zero, and countable additivity gives
$\mu(\CC^*)=0$.  The half-open annuli assign every boundary circle to
exactly one set, so this argument applies equally to atomic and diffuse
measures and to nonreal $\lambda$.

Thus the nonempty support of $\mu$ consists of one or two of the fixed
points of $T$.  Invariance under $A$ and $B$ preserves this support as a
set: the inverse image of an open neighborhood of the image of a
support point has positive measure.  Both projective transformations
have order three.  Their induced permutations on a set of at most two
points are therefore trivial.  Each point of the support would be
fixed by both $A$ and $B$, and hence would be a common eigenline of the
two matrices. Lemma~\ref{lem:matrix-data} excludes such a line.
\end{proof}

The preceding results rule out exactly the object supplied by
Proposition~\ref{prop:limit}.

\begin{proposition}
\label{prop:obstruction}
For the bundle $G$ of Proposition~\ref{prop:construction}, every finite
nonnegative plurisubharmonic function on the total space of $F=G^\vee$
whose restriction to each fiber is a Hermitian quadratic form is
identically zero.
\end{proposition}

\begin{proof}
If such a function were nonzero, Lemma~\ref{lem:ruling-reduction}
would give a nonzero squared seminorm on $\OO_{\PP^1}(-b)$ invariant
under the lifts of $A$ and $B$. Lemma~\ref{lem:seminorm-curvature}
would then produce a finite positive measure on $\PP^1$ of mass
$2\pi b>0$, invariant under both projective transformations.
This contradicts Lemma~\ref{lem:no-invariant-measure}.
\end{proof}
